% Generated draft: requires mathematical and source review.
\documentclass[11pt]{article}
\usepackage[margin=25mm]{geometry}
\usepackage{fontspec}
\setmainfont{Latin Modern Roman}
\usepackage{amsmath,amssymb,mathtools,unicode-math}
\setmathfont{Latin Modern Math}
\usepackage{enumitem,needspace,xurl,hyperref}
\setlength{\parindent}{0pt}
\setlength{\parskip}{5pt}
\setlength{\emergencystretch}{3em}
\newcommand{\N}{\mathbb{N}}
\newcommand{\Z}{\mathbb{Z}}
\newcommand{\Q}{\mathbb{Q}}
\newcommand{\R}{\mathbb{R}}
\newcommand{\C}{\mathbb{C}}
\newcommand{\F}{\mathbb{F}}
\newcommand{\A}{\mathbb{A}}
\newcommand{\PP}{\mathbb{P}}
\newcommand{\E}{\mathbb{E}}
\newcommand{\eps}{\varepsilon}
\newcommand{\dd}{\,\mathrm{d}}
\begin{document}
\section*{UnsolvedMath ID: 701154}
\textbf{Source problem code:} catalogue entry
\subsection{Definitions and mathematical statement}
A variety of groups is a class consisting of all groups satisfying a specified collection of group identities. A subvariety is a subclass that is itself a variety. A variety is generated by its finite members if it is the smallest variety containing those members; this is the interpretation used here for the source phrase ``defined by its finite groups.'' For a class \(\mathfrak{X}\) and an integer \(n\geq 1\), define
\[
g_{\mathfrak{X}}(n)=\#\{[G]:G\in\mathfrak{X},\ |G|=n\},
\]
where \([G]\) denotes the isomorphism class of \(G\).

The general problem asks for good upper bounds for
\[
\frac{g_{\mathfrak{V}}(n)}{g_{\mathfrak{U}}(n)},
\qquad g_{\mathfrak{U}}(n)>0,
\]
when \(\mathfrak{V}\) is a variety generated by its finite members and \(\mathfrak{U}\subseteq\mathfrak{V}\) is a subvariety. The source does not specify a quantitative meaning for ``good'' or a proposed general bound.

A group is a \(p\)-group, for a prime \(p\), if each of its elements has finite order a power of \(p\). A variety is locally finite if every finitely generated group in it is finite. A non-abelian variety means a variety containing at least one non-abelian group.

\textbf{Conjecture attributed to P. M. Neumann.} Let \(p\) be a prime, let \(\mathfrak{V}\) be a locally finite variety of \(p\)-groups, and let \(\mathfrak{U}\subseteq\mathfrak{V}\) be a non-abelian subvariety. The source conjectures
\[
\frac{g_{\mathfrak{V}}(p^m)}{g_{\mathfrak{U}}(p^m)}<p^{O(m^2)}.
\]
Interpreting this as an asymptotic upper bound for each fixed triple \((p,\mathfrak{V},\mathfrak{U})\), its precise formulation is: there exist a real constant \(C>0\) and an integer \(m_0\geq 1\) such that, for every integer \(m\geq m_0\) with \(g_{\mathfrak{U}}(p^m)>0\),
\[
\frac{g_{\mathfrak{V}}(p^m)}{g_{\mathfrak{U}}(p^m)}<p^{Cm^2}.
\]
The source does not specify whether constants should be uniform in the prime or in either variety; the formulation above makes no such uniformity assertion.

For context, let \(a_p(m)\) be the number of isomorphism classes of all groups of order \(p^m\), and define
\[
\varepsilon_p(m)=\log_p a_p(m)-\frac{2}{27}m^3.
\]
The source also describes the Sims Conjecture as the assertion \(\varepsilon_p(m)=O(m^2)\), equivalently that for each fixed prime \(p\) there exist \(K>0\) and \(M\geq1\) such that \(|\varepsilon_p(m)|\leq Km^2\) for every integer \(m\geq M\). Uniformity in \(p\) is not specified. The source suggests that the variety-counting conjecture might provide an approach to this conjecture; it does not supply an implication or proof.
\subsection{Short English statement}
Count finite groups up to isomorphism within a variety. Neumann asks for upper bounds on the ratio of these counts for a variety and a subvariety, wherever the denominator is positive. His specific conjecture predicts an upper bound of the form \(p^{Cm^2}\) for locally finite varieties of \(p\)-groups and non-abelian subvarieties. The source suggests a possible connection to the Sims Conjecture on the error term in counting all groups of order \(p^m\).
\subsection{Sources}
\begin{itemize}
\item Local UnsolvedMath registry, supplied complete source record with permanent ID \(701154\), exact source problem code \texttt{catalogue entry}, and title ``Kourovka Notebook Problem 15.66.'' The record labels the problem open; this is supplied registry status, not an independently verified mathematical assessment. The record was created and updated on 2026-10-16 and identifies upstream revision \texttt{d371ee917ac9647fbb06cd6fc529deec6f3dda82}.
\item P. M. Neumann, proposer, 2002. The supplied attribution identifies \emph{The Kourovka Notebook (Unsolved Problems in Group Theory)}, problem 15.66, 15th issue (2002), as reproduced in the 21st issue edition (2026), arXiv:1401.0300v48. The registry describes the collection as edited by E. I. Khukhro and V. D. Mazurov. These bibliographic details are supplied and have not been independently checked.
\item G. Venkataraman, \emph{Quart. J. Math. Oxford (2)} \textbf{48}, no. 189 (1997), 107--125. The source reports progress for varieties generated by finite groups all of whose Sylow subgroups are abelian. No article title or precise result is supplied.
\item The Sims Conjecture and its proposed connection to the main conjecture are mentioned in the supplied problem statement. No separate bibliographic reference for that conjecture is supplied.
\end{itemize}
\end{document}
